\documentclass{article}
\input{../../lib.sty}


\title{\texttt{fn zcdp\_delta}}
\author{\MichaelShoemateAuthor}

\begin{document}
\maketitle

Proves soundness of \texttt{fn zcdp\_delta}.
This proof is an adaptation of \href{https://arxiv.org/pdf/2004.00010.pdf#subsection.2.3}{subsection 2.3} of \cite{CKS20}.
\section{Bound Derivation}

\begin{definition}
\label{plrv}
(Privacy Loss Random Variable). 
Let $M: \mathcal{X}^n \rightarrow \mathbb{Y}$ be a randomized algorithm.
Let $x, x' \in \mathcal{X}^n$ be neighboring inputs.
Define $f: \mathcal{Y} \rightarrow \mathbb{R}$ by $f(y) = log \left( \frac{\mathbb{P}[M(x)=y]}{\mathbb{P}[M(x')=y]} \right)$.
Let $Z = f(M(x))$, the privacy loss random variable, denoted $Z \leftarrow PrivLoss(M(x) || M(x'))$.
\end{definition}

\begin{lemma}
\label{delta-bound}
\cite{CKS20} Let $\epsilon, \delta \geq 0$. Let M: $\mathcal{X}^n \rightarrow \mathcal{Y}$ be a randomized algorithm. Then M satisfies $(\epsilon, \delta)$-differential privacy if and only if

\begin{align}
    \delta &\geq \underset{Z \leftarrow PrivLoss(M(x)|| M(x'))}{\mathbb{E}} [max(0, 1 - e^{\epsilon - Z})] \\
    % &= \underset{Z \leftarrow PrivLoss(M(x)|| M(x'))}{\mathbb{P}} [Z > \epsilon] - e^\epsilon \underset{Z' \leftarrow PrivLoss(M(x')|| M(x))}{\mathbb{P}} [-Z' > \epsilon] \\
    % &= \int_\epsilon^\infty e^{\epsilon - z} \underset{Z \leftarrow PrivLoss(M(x)|| M(x'))}{\mathbb{P}} [Z > z]dz
\end{align}
for all $x, x' \in \mathcal{X}^n$ differing on a single element.

\begin{proof}
Fix neighboring inputs $x, x' \in \mathcal{X}^n$. 
Let $f: \mathcal{Y} \rightarrow \mathbb{R}$ be as in \ref{plrv}.
For notational simplicity, let $Y = M(x)$, $Y' = M(x')$, $Z = f(Y)$ and $Z' = -f(Y')$.
This is equivalent to $Z \leftarrow PrivLoss(M(x) || M(x'))$.
Our first goal is to prove that

\begin{equation}
    \sup\limits_{E \subset \mathcal{Y}} \mathbb{P}[Y \in E] - e^\epsilon \mathbb{P}[Y' \in E] = \mathbb{E}[max\{0, 1 - e^{\epsilon - Z}\}].
\end{equation}

For any $E \subset \mathcal{Y}$, we have
\begin{equation}
    \mathbb{P}[Y' \in E] = \mathbb{E}[\mathbb{I}[Y' \in E]] = \mathbb{E}[\mathbb{I}[Y \in E] e^{-f(Y)}].
\end{equation}
This is because $e^{-f(y)} = \frac{\mathbb{P}[Y'=y]}{\mathbb{P}[Y=y]}$.

Thus, for all $E \subset \mathcal{Y}$, we have 
\begin{equation}
    \mathbb{P}[Y \in E] - e^\epsilon \mathbb{P}[Y' \in E] = \mathbb{E} \left[ \mathbb{I}[Y \in E] (1 - e^{\epsilon - f(Y)}) \right]
\end{equation}

Now it is easy to identify the worst event as $E = \{y \in \mathcal{Y} : 1 - e^{\epsilon - f(y)} > 0\}$. Thus
\begin{equation}
    \sup\limits_{E \subset Y} \mathbb{P}[Y \in E] - e^{\epsilon} \mathbb{P}[Y' \in E] = \mathbb{E} \left[ \mathbb{I}[1 - e^{\epsilon - f(Y)} > 0] (1 - e^{\epsilon - f(Y)}) \right] = \mathbb{E}[max\{ 0, 1 - e^{\epsilon - Z} \}]
\end{equation}

% Alternatively, since the worst event is equivalently $E = \{y \in \mathcal{Y} : f(y) > \epsilon \}$, we have

% \begin{equation}
%     \sup\limits_{E \subset Y} \mathbb{P}[Y \in E] - e^{\epsilon} \mathbb{P}[Y' \in E] =
%     \mathbb{P} [f(Y) > \epsilon] - e^\epsilon \mathbb{P} [f(Y') > \epsilon] =
%     \mathbb{P} [Z > \epsilon] - e^\epsilon \mathbb{P}[-Z' > \epsilon]
% \end{equation}

% It only remains to show that
% \begin{equation}
%     \mathbb{E} [max \{0, 1 - e^{\epsilon - Z}\}] = \int_\epsilon^\infty e^{\epsilon - z} \mathbb{P}[Z > z] dz.
% \end{equation}

% This follows from integration by parts: Let $u(z) = \mathbb{P}[Z > z]$ and $v(z) = 1 - e^{\epsilon - z}$ and $w(z) = u(z)v(z)$. Then
% \begin{align}
%     \mathbb{E}[max\{ 0, 1 - e^{\epsilon - Z} \}] &= \int_\epsilon^\infty v(z) u'(z) dz = \int_\epsilon^\infty (w'(z) - v'(z) u(z)) dz
%     &= \lim\limits_{z \rightarrow \infty} w(z) - w(\epsilon) + \int_\epsilon^\infty e^{\epsilon - z} \mathbb{P} [Z > z] dz
% \end{align}

% Now $w(\epsilon) = u(\epsilon) (1 - e^{\epsilon - \epsilon} = 0$ and $0 \leq \lim_{z \rightarrow \infty} w(z) \leq \lim_{z \rightarrow \infty} \mathbb{P}[Z > z] = 0$, as required.

\end{proof}

\end{lemma}

\begin{theorem}
\label{renyidp-approxdp-delta}
\cite{CKS20} Let $M: \mathcal{X}^n \rightarrow \mathcal{Y}$ be a randomized algorithm. Let $\alpha \in (1, \infty)$ and $\epsilon \geq 0$. Suppose $D_\alpha(M(x) || M(x')) \leq \tau$ for all $x, x' \in \mathcal{X}^n$ differing in a single entry.\footnote{This is the definition of $(\alpha, \tau)$-R\'enyi differential privacy.} Then M is $(\epsilon, \delta)$-differentially private for 

\begin{equation}
    \delta = \frac{e^{(\alpha - 1) (\tau - \epsilon)}}{\alpha - 1} \left(1 - \frac{1}{\alpha}\right)^\alpha
\end{equation}
\end{theorem}

\begin{proof}
Fix neighboring $x, x' \in  \mathcal{X}^n$ and let $Z \leftarrow PrivLoss(M(x) || M(x'))$. We have
\begin{equation}
    \mathbb{E}[e^{(\alpha - 1) Z}] = e^{(\alpha - 1)D_\alpha(M(x) || M(x'))} \leq e^{(\alpha - 1)\tau}
\end{equation}

By \ref{delta-bound}, our goal is to prove that $\delta \geq \mathbb{E}[max\{0, 1 - e^{\epsilon - Z} \}]$. 
Our approach is to pick $c > 0$ such that $max\{0, 1 - e^{\epsilon - Z} \} \leq c e^{(\alpha - 1) z}$ for all $z \in \mathbb{R}$. Then
\begin{equation}
    \mathbb{E}[max\{0, 1 - e^{\epsilon - Z} \}] \leq \mathbb{E}[c e^{(\alpha - 1) z}] \leq c e^{(\alpha - 1) \tau}.
\end{equation}

We identify the smallest possible value of $c$:
\begin{equation}
    c = \sup\limits_{z \in \mathbb{R}} \frac{max\{0, 1 - e^{\epsilon - z} \}}{e^{(\alpha -1)z}} 
    = \sup\limits_{z \in \mathbb{R}} e^{z - \alpha z} - e^{\epsilon - \alpha z} 
    = \sup\limits_{z \in \mathbb{R}}f(z)
\end{equation}
where $f(z) = e^{z - \alpha z} - e^{\epsilon - \alpha z}$. We have 

\begin{equation}
    f'(z) = e^{z - \alpha z} (1 - \alpha) - e^{\epsilon - \alpha z}(-\alpha) 
    = e^{-\alpha z} (\alpha e^\epsilon - (\alpha - 1) e^z)
\end{equation}

Clearly $f'(z) = 0 \Longleftrightarrow e^z = \frac{\alpha}{\alpha - 1}e^{\epsilon} \Longleftrightarrow z = \epsilon - log(1 - 1/\alpha)$. Thus
\begin{align}
    c &= f(\epsilon - log(1 - 1 / \alpha)) \\
    &= \left(\frac{\alpha}{\alpha - 1} e^\epsilon \right)^{1 - \alpha} - e^{\epsilon} \left(\frac{\alpha}{\alpha - 1} e^\epsilon \right)^{- \alpha} \\
    &= \left(\frac{\alpha}{\alpha - 1} e^\epsilon  - e^\epsilon \right)  \left(\frac{\alpha}{\alpha - 1} e^{-\epsilon} \right)^{\alpha} \\
    &= \frac{e^\epsilon}{\alpha - 1} \left(1 - \frac{1}{\alpha} \right)^\alpha  e^{-\alpha \epsilon}.
\end{align}

Thus
\begin{equation}
    \mathbb{E}[max\{0, 1 - e^{\epsilon - Z}\}] \leq \frac{e^\epsilon}{\alpha - 1}\left(1  - \frac{1}{\alpha} \right) ^\alpha e^{-\alpha \epsilon} e^{(\alpha -1)\tau}
    = \frac{e^{(\alpha - 1)(\tau - \epsilon)}}{\alpha - 1} \left( 1 - \frac{1}{\alpha} \right)^\alpha
    = \delta
\end{equation}

\end{proof}

\begin{theorem}
\label{asoodeh-rdp-approxdp}
\cite[Lemma 1, Equation (25)]{ALCKS21}
Let $M$ satisfy $(\alpha,\gamma)$-RDP for $\alpha>1$. If
$\epsilon>\gamma$, then $M$ is $(\epsilon,\delta)$-differentially
private for
\begin{equation}
    \delta =
    \frac{e^{(\alpha-1)\gamma}-1}
         {\alpha\left(e^{(\alpha-1)\epsilon}-1\right)}.
\end{equation}
\end{theorem}

\begin{proof}
The second term of \cite[Lemma 1, Equation (25)]{ALCKS21} states that,
when $0<\alpha\delta<1$, the smallest guaranteed approximate-DP
epsilon is at most
\begin{equation}
    \frac{1}{\alpha-1}
    \log\left(
        \frac{e^{(\alpha-1)\gamma}-1}{\alpha\delta}+1
    \right).
\end{equation}
Substitution of the stated $\delta$ makes this expression equal to
$\epsilon$. Moreover, $\epsilon>\gamma$ implies
$e^{(\alpha-1)\epsilon}-1>e^{(\alpha-1)\gamma}-1$, and therefore
$0<\alpha\delta<1$, satisfying the lemma's condition.
\end{proof}

\begin{corollary}
\label{renyidp-approxdp-epsilon}
\cite{CKS20} Let $M: \mathcal{X}^n \rightarrow \mathcal{Y}$ be a randomized algorithm. Let $\alpha \in (1, \infty)$ and $\epsilon \geq 0$. Suppose $D_\alpha(M(x) || M(x')) \leq \tau$ for all $x, x' \in \mathcal{X}^n$ differing in a single entry. Then M is $(\epsilon, \delta)$-differentially private for 

\begin{equation}
    \epsilon = \tau + \frac{\ln(1 / \delta) + (\alpha - 1) \ln(1 - 1/\alpha) - \ln(\alpha)}{\alpha - 1}
\end{equation}
\end{corollary}

\begin{proof}
This follows by rearranging \ref{renyidp-approxdp-delta}.
\end{proof}

\begin{corollary}
\label{zcdp-approxdp-epsilon}
Let $M: \mathcal{X}^n \rightarrow \mathcal{Y}$ be a randomized algorithm satisfying $\rho$-concentrated differential privacy. Then M is $(\epsilon, \delta)$-differentially private for any $0 < \delta \leq 1$ and 
\begin{equation}
    \epsilon = \inf\limits_{\alpha \in (1, \infty)} \alpha \rho + \frac{\ln(1 / \delta) + (\alpha - 1) \ln(1 - 1/\alpha) - \ln(\alpha)}{\alpha - 1}
\end{equation}
\end{corollary}

\begin{proof}
This follows from \ref{renyidp-approxdp-epsilon} by taking the infimum over all divergence parameters $\alpha$.
\end{proof}

\section{Certified inverse}

For a fixed order $\alpha$, the expression in
\ref{renyidp-approxdp-epsilon} is the inverse of the branch-zero
log-delta bound. The directed interval implementation evaluates this
expression with $\gamma = \alpha\rho$ and returns an upper endpoint, so
that endpoint is a sufficient epsilon for the requested delta.

For the second branch, the theorem requires $\epsilon > \gamma$ and its
inverse additionally requires $\alpha\delta < 1$. Solving the bound in
\ref{asoodeh-rdp-approxdp} gives
\begin{equation}
    \epsilon = \frac{1}{\alpha-1}
    \operatorname{softplus}\left(\log(\exp((\alpha-1)\gamma)-1)
    -\log(\alpha)-\log(\delta)\right).
\end{equation}
The implementation represents a branch that cannot establish these strict
domains as `None`; numerical failures remain errors. Thus every certified
candidate is independently a valid upper bound, and the minimum of the
candidates is valid as well.

Native floating-point arithmetic is used only to select promising orders.
The final bracket endpoints and optimizer argument are recomputed with
outward-rounded `DInterval` arithmetic. Search accuracy therefore affects
only tightness, not soundness. The special cases $\rho=0$, $\rho=\infty$,
$\delta=0$, and $\delta=1$ return the corresponding trivial valid bounds
before finite interval evaluation.

\begin{theorem}
For any valid $\rho$ and $\delta$, a successful return from
\texttt{zcdp\_epsilon} is a sufficient epsilon for the requested delta.
\end{theorem}
\begin{proof}
The fixed-order inverse formulas above are sufficient by their respective
forward theorems. The optimizer only chooses the finite candidate orders;
each candidate is certified independently with outward-rounded interval
arithmetic. The upper endpoint is therefore conservative, and the minimum
of independently valid bounds remains valid. The explicit special-case
returns are valid by the trivial privacy bounds. Errors are propagated, so
no uncertified arithmetic result is accepted.
\end{proof}

\section{Hoare Triple}
\subsection*{Precondition}
\subsubsection*{Compiler-verified}
None

\subsubsection*{Caller-verified}
None

\subsection*{Pseudocode}
\lstinputlisting[language=Python,firstline=2]{./pseudocode/cdp_delta.py}


\subsection*{Postcondition}

\begin{theorem}
    For any possible setting of $\rho$ and $\epsilon$, $\texttt{zcdp\_delta}$ either returns an error,
    or a $\delta$ such that any $\rho$-differentially private measurement is also $(\epsilon, \delta)$-differentially private.
\end{theorem}
\begin{proof}
The validation and special-case branches return an error or a trivial valid
bound. Otherwise, ordinary floating-point optimization only selects a finite
set $A$ of candidate orders strictly greater than one. Its accuracy affects
tightness but not soundness.

For each candidate, the code evaluates the logarithm of either the bound from
\ref{renyidp-approxdp-delta}, or, when $\epsilon>\alpha\rho$, the bound from
\ref{asoodeh-rdp-approxdp}, with the zCDP guarantee
$\gamma=\alpha\rho$. These evaluations use \texttt{DInterval}, whose
operations round outwards, so each upper endpoint is at least the exact
logarithm of the corresponding valid bound. The minimum remains valid because
every candidate bounds the same privacy quantity from above.

Finally, \texttt{InfExp} rounds exponentiation upwards and returns the smallest
positive \texttt{f64} on underflow. Any interval-evaluation failure is
returned as an error rather than being used as a privacy result. Thus every
successful returned \texttt{delta} satisfies the postcondition.
\end{proof}


\bibliographystyle{alpha}
\bibliography{mod}

\end{document}
